\documentclass[10pt]{article}

\usepackage[utf8]{inputenc}

\usepackage[T1]{fontenc}

\usepackage{amsmath}

\usepackage{amsfonts}

\usepackage{amssymb}

\usepackage[version=4]{mhchem}

\usepackage{stmaryrd}

\usepackage{bbold}



\begin{document}

КАТАСТРОФ ТЕОРИЯ - совокупность приложений теории особенностей дифференцируемых (гладких) отображений Х. Уитни (H. Whitney) и теории бифуркаций А. Пуанкаpe (Н. Poincaré) и А. А. Андронова. Название введено Р. Томом (R. Thom) в 1972. К. т. применяется к геометрич. и физич. оптике, гидродинамике, устойчивости кораблей, а также к исследованию биений сердца, эмбриологии, социологии, лингвистике, экспериментальной психологии, экономике, геологии, теории элементарных частиц и моделированию деятельности мозга и психич. расстройств и т. п. Поскольку гладкие отображения встречаются повсеместно, неудивительно, что повсеместно встречаются и их особенности. Когда явление описывается гладким отображением и нет причин для нетипичности (напр., симметрий), применение теории особенностей оправдано и полезно (в оптике, теории упругости и др.), тогда как в нек-рых из описанных P. Томом и Э. Зиманом (E. Zeeman) приложений сомнительно уже существование изучаемого отображения (в биологии, лингвистике, социологии).



Теория особенностей обобщает исследование экстремумов функций на случай нескольких функций любого числа переменных. К р и т и чес к о й точ ко й функции у называется точка, в к-рой все первые частные производные равны нулю: $\partial y / \partial x_{i}=0$; критич. точка называется не выро ж д е н н о й, если матрица $\left\|\partial^{2} y / \partial x_{i} \partial x_{j}\right\|$ невырождена, то есть ее определитель отличен от нуля. У типичной функции все критич. точки невырождены. Любая гладкая функция в окрестности каждой невырожденной критич. точки приводится к одной из так наз. нор мальных форм Морса $y= \pm x_{1}^{2} \pm \ldots \pm x_{n}^{2}+C$, гладкой заменой независимых переменных. Эти невырожденные особенности устойчивы: напр., всякая функция, достаточно близкая к $y=x^{2}$ (с производными), имеет в подходящей точке вблизи нуля подобную же особенность (невырожденную точку минимума). Все более сложные особенности неустойчивы. Напр., вырожденная критич. точка функции $y=x^{3}$ в нуле распадается на две при возмущении, превращающем $x^{3}$ в $x^{3}-\varepsilon x$.



Типичные отображения поверхности на плоскость $\left(\mathbb{R}^{2} \rightarrow \mathbb{R}^{2}\right.$ ) также имеют лишь устойчивые особенности, а именно, складку $\left(y_{1}=x_{1}^{2}, y_{2}=x_{2}\right)$ либо с борку Уитни $\left(y_{1}=x_{1}^{3}+x_{1} x_{2}, y_{2}=x_{2}\right)$. Сборка есть особенность проектирования поверхности $y_{1}=x_{1}^{3}+x_{1} x_{2}$ из пространства $\left(x_{1}, x_{2}, y_{1}\right)$ на плоскость $\left(y_{1}, x_{2}\right)$ (рис. 1). Списки типичных особенностей отображений $\mathbb{R}^{3} \rightarrow \mathbb{R}^{3}$ и $\mathbb{R}^{4} \rightarrow \mathbb{R}^{4}$ таковы: 1) $y_{1}=x_{1}^{2}$, $y_{i}=x_{i}, i>1$; 2) $y_{1}=x_{1}^{3}+x_{1} x_{2}, y_{i}=x_{i}, i>1$; 3) $y_{1}=x_{1}^{4}+x_{1}^{2} x_{2}+$ $+x_{1} x_{3}, \quad y_{i}=x_{i}, \quad i>1 ;$ 4) $y_{1}=x_{1}^{2} \pm x_{2}^{2}+x_{1} x_{3}+x_{2} x_{4}, \quad y_{2}=x_{1} x_{2}$, $y_{3}=x_{3}, y_{4}=x_{4}$. Отображение $\mathbb{R}^{2} \rightarrow \mathbb{R}^{3}$ обычно имеет особенностями лишь зонтики Уитни - Кэли (рис. 2 ; $y_{1}=x_{1}^{2}, y_{2}=x_{1} x_{2}, y_{3}=x_{2}$ ). При переходе к высшим размерностям списки типичных особенностей растут и даже становятся континуальными (напр., не всякое отображение $\mathbb{R}^{n} \rightarrow \mathbb{R}^{n}$ при $n>8$ аппроксимируется устойчивым). Число классов топологически различных особенностей остается конечным при любых размерностях.



В теории бифуркаций рассматривается динамич. система, описываемая уравнением $\dot{x}=\theta(x, \varepsilon)$, с заданным векторным полем $\theta$ в $n$-мерном фазовом пространстве $\{x\}$. Поле зависит от $k$-мерного параметра $\varepsilon$. Множество состояний равновесия определяет в $(n+k)$-мерном пространстве $\{x, \varepsilon\} k$-мерную поверхность $\theta(x, \varepsilon)=0$. В типичном случае эта поверхность гладкая, но ее проекция на пространство «управляющих параметров» $\{\varepsilon\}$ может иметь особенности. Если рассматривать значения $\{\varepsilon\}$ как функции на поверхности



238 КАТАCTPOФ состояний равновесия, то точки, в к-рых якобиан этих функций равен 0, называются б и фур ка ци о н ны м и, а значения функций в этих точках - бифуркационными значениями параметров $\varepsilon$. При подходе управляющих параметров к бифуркационным значениям положения равновесия «бифуркацируют» (рождаются или умирают). Знание геометрии типичных особенностей позволяет описывать происходящие при этом явления, напр. скачкообразный переход системы к далекому состоянию равновесия при плавном изменении параметров. Такие скачки способны разрушить систему (механическую, упругую, электрическую, биологическую, химическую и т. п.), откуда и название «К. т.».



Наибольший успех достигнут в приложениях К. т. к оптике, где даже типичные особенности каустик и перестройки волновых фронтов в трехмерном пространстве не были известны. Рассмотрим возмущение (свет, звук, ударную волну, эпидемию и др.), распространяюшееся с единичной скоростью из области, ограниченной гладким фронтом. Чтобы построить фронт через время $t$, нужно отложить отрезок длины $t$ на каждом луче нормали. Через нек-рое время на движущемся фронте появляются особенности в точках каустики (огибающей семейства лучей) исходного фронта. Напр., при распространении возмущения внутрь эллипса на плоскости особенности фронта скользят по каустике, имеющей 4 точки возврата (рис. 3). Эти особенности устойчивы (не исчезают при малой деформации исходного фронта). Типичные особенности фронтов в трехмерном пространстве - это самопересечения, ребра возврата (нормальная форма $x^{2}=y^{3}$ ) и лас точкины хвосты [рис. 4; эта поверхность образована точками $(a, b, c)$, для к-рых многочлен $x^{4}+a x^{2}+b x+c$ имеет кратный корень]. Каустики в трехмерном пространстве имеют особенности еще двух видов (п и рам ида и кошелек; рис. 5).



Почти все особенности волновых фронтов (или Лежандра преобразований) можно описать как множества бифуркационных значений параметра $\mu$, при к-рых возникают особенности отображения $(x, \mu) \rightarrow \mu$ гиперповерхности $F(x, \mu)=0$ в пространство $\mu$, где $F$ - типичное семейство гладких функций вектора $x$ и векторного параметра $\mu$. Типичные особенности каустик (или градиентных отображений $x \rightarrow \partial S / \partial x$, или отображений Гаусса, сопоставляющих точке поверхности направление нормали) можно описать как множества бифуркационных значений параметра $\mu$, при к-рых функция $F(x, \mu)$ переменной $x$ имеет вырожденную критич. точку. Ласточкин хвост, пирамида и кошелек получаются при



\[

\begin{gathered}

F=x^{5}+\mu_{1} x^{2}+\mu_{2} x^{2}+\mu_{3} x ; \\

F=x_{1}^{2} x_{2} \pm x_{2}^{3}+\mu_{1} x_{2}^{2}+\mu_{2} x_{2}+\mu_{3} x_{1} .

\end{gathered}

\]



Особенностям каустик и фронтов геометрич. оптики соответствуют в волновой теории особенности асимптотик осциллируюших интегралов в методе стационарной фазы или многомерном методе перевала при слиянии нескольких стационарных точек. По порядку величины интеграл при подходе к точке каустики возрастает в $\lambda^{-\nu}$ раз, где $\lambda-$ длина волны, а показатель $v$ равен $1 / 6$ для общей точки каустики $\left(A_{2}\right.$, особенность Эйри); $1 / 4$ для общей точки ребра возврата ( $A_{3}$, особенность Пирси); $3 / 10$ для ласточкина хвоста (особенность $A_{4}$ ); 1/3 для кошелька и пирамиды (особенности $D_{4}$ ). Эти особенности связаны с простыми группами Ли $A_{k} \sim S U(k+1), D_{k} \sim O(2 k)$, а также с правильными многогранниками [конечными подгруппами группы $S U(2)]$. Показатель $v$ определяет интенсивность света вблизи каустики и ее особенностей, разрушение среды интенсивной волной, скопление частиц при движении пылевидной среды с потенциальным полем скоростей (с иным значением v) и т. п. Универсальность геометрии бифуркационных диаграмм позволяет использовать их для одновременного моделирования многих различных по своему физич. смыслу явлений.





\end{document}